%% =====================================================================
\chapter{Authentication and Sessions}\label{ch:auth}

This chapter covers who a user is (the \code{User} model), how they prove it
(password plus emailed OTP), and how the proof is carried between requests
(access and refresh tokens).

\section{The User model: \file{src\codesl{}models\codesl{}User\codedot{}js}}

One collection holds all four roles. Fields that only make sense for one role
live in a sub-document (\code{patientProfile}, \code{doctorProfile},
\code{agentProfile}).

\begin{lstlisting}[language=JavaScript,caption={User schema, key fields (\file{src/models/User.js})}]
const ROLES = ['patient', 'doctor', 'admin', 'insurance_agent'];
// Doctors and insurance agents must be approved by an admin before they can
// touch clinical or financial data. Patients and admins are approved on creation.
const ROLES_NEEDING_APPROVAL = ['doctor', 'insurance_agent'];

const userSchema = new mongoose.Schema({
  email: { type: String, required: true, unique: true, lowercase: true, trim: true },
  fullName: { type: String, trim: true, required: true },
  passwordHash: { type: String, required: true },
  role: { type: String, enum: ROLES, required: true, default: 'patient' },
  // ... phone, dateOfBirth, gender, address, patientProfile, doctorProfile, agentProfile
  approvalStatus: { type: String, enum: ['pending', 'approved', 'rejected'], default: 'approved' },
  approvalNote: { type: String, default: '' },

  // Identity of this user inside the hash-chain ledger (was `walletAddress`).
  ledgerAddress: { type: String, index: true },

  isVerified: { type: Boolean, default: false }, // proven control of the email (first OTP login)
  isActive: { type: Boolean, default: true },
  lastLogin: { type: Date, default: null },
  failedLoginAttempts: { type: Number, default: 0 },
  accountLockedUntil: { type: Date, default: null },
  refreshTokens: {
    type: [String],
    default: [],
    validate: { validator: (tokens) => tokens.length <= 5, message: 'A user can have at most 5 refresh tokens' },
  },
}, { timestamps: true });
\end{lstlisting}

A doctor's \code{doctorProfile} also holds \code{availability}: working days
(0 = Sunday), start and end hour, and slot length. Appointments are generated
from it (Chapter~\ref{ch:appointments}).

\begin{lstlisting}[language=JavaScript,caption={Password hashing and ledger address (\file{src/models/User.js})}]
userSchema.pre('save', async function hashPasswordAndAssignLedgerAddress(next) {
  try {
    if (!this.ledgerAddress) {
      this.ledgerAddress = `mc_${sha256(this._id.toString()).slice(0, 40)}`;
    }
    if (this.isModified('passwordHash')) {
      const salt = await bcrypt.genSalt(10);
      this.passwordHash = await bcrypt.hash(this.passwordHash, salt);
    }
    next();
  } catch (error) {
    next(error);
  }
});
\end{lstlisting}

\begin{howto}
Code that creates or changes a user puts the \emph{plain} password into
\code{passwordHash}; this hook replaces it with a bcrypt hash before anything is
written. \code{isModified} ensures an already-hashed password is not hashed
again on an unrelated save. bcrypt is deliberately slow (cost 10, about 2\textsuperscript{10}
rounds), so a stolen database cannot be brute-forced quickly. The
\code{ledgerAddress} is a stable pseudonymous identifier derived from the user
id, the successor of the paper's blockchain wallet address.
\end{howto}

\begin{lstlisting}[language=JavaScript,caption={Never leaking secrets in JSON (\file{src/models/User.js})}]
userSchema.methods.toJSON = function toJSON() {
  const obj = this.toObject();
  delete obj.passwordHash;
  delete obj.failedLoginAttempts;
  delete obj.accountLockedUntil;
  delete obj.refreshTokens;
  delete obj.__v;
  // Hide the other roles' (empty) profile sections. When a partial populate did
  // not select `role`, keep whatever profile fields were explicitly selected.
  if (obj.role) {
    if (obj.role !== 'patient') delete obj.patientProfile;
    if (obj.role !== 'doctor') delete obj.doctorProfile;
    if (obj.role !== 'insurance_agent') delete obj.agentProfile;
  }
  return obj;
};
\end{lstlisting}

Express calls \code{toJSON} whenever a user is sent with \code{res.json}, so no
route can accidentally return a password hash. The \code{if (obj.role)} guard
fixes a real bug found during testing: when an appointment populated only
\code{doctorProfile.specialization} (without \code{role}), the doctor's profile
was stripped and the patient dashboard crashed.

\section{Registration}

\begin{lstlisting}[language=JavaScript,caption={register, core (\file{src/modules/auth/auth.controller.js})}]
async function register(req, res) {
  const { email, password, fullName, role, phone, dateOfBirth, gender, address } = req.body;

  if (await User.exists({ email })) {
    throw conflict('An account with this email already exists', 'EMAIL_TAKEN');
  }

  const needsApproval = User.ROLES_NEEDING_APPROVAL.includes(role);
  const user = new User({
    email,
    fullName,
    passwordHash: password, // hashed by the pre-save hook
    role,
    // ...
    approvalStatus: needsApproval ? 'pending' : 'approved',
  });
  // ... role-specific profile fields
  await user.save();
  await ledgerService.append({
    chain: 'records', entityType: 'user', entityId: user._id, action: 'REGISTER',
    actorId: user._id, actorRole: role, payload: { email: user.email, role, ledgerAddress: user.ledgerAddress },
  });

  if (needsApproval) {
    await notificationService.notifyRole('admin', {
      type: 'warning',
      title: 'New account awaiting verification',
      // ...
      page: 'verification',
    });
  }

  emailService.sendWelcome(email, fullName, role, needsApproval); // fire-and-forget by design
\end{lstlisting}

\begin{howto}
The Joi schema on the route only allows \code{role} to be \code{patient},
\code{doctor} or \code{insurance\_agent}; nobody can register as an admin. Doctors
must supply a specialization and licence number, agents a company. A doctor or
agent account starts \code{pending}: they can log in and edit their profile, but
every clinical or claims endpoint answers 403 \code{APPROVAL\_PENDING} until an
admin approves them (Section~\ref{sec:policy}). The welcome email is not awaited:
if the mail server is slow or down, registration still succeeds.
\end{howto}

\section{Two-step login}

\begin{figure}[H]
\centering
\begin{tikzpicture}[node distance=0.5cm]
  \node[box,minimum width=3cm] (b1) {Browser};
  \node[box,minimum width=3cm,right=4.4cm of b1] (a1) {API};
  \node[box,minimum width=2.4cm,right=3cm of a1] (r1) {Redis};
  \draw[arr] ([yshift=2mm]b1.east) -- node[lbl,above]{1. email + password} ([yshift=2mm]a1.west);
  \draw[arr] ([yshift=2mm]a1.east) -- node[lbl,above]{store SHA-256(code)} ([yshift=2mm]r1.west);
  \node[below=0.2cm of a1,font=\scriptsize] (mail) {code emailed (SMTP)};
  \node[box,minimum width=3cm,below=1.6cm of b1] (b2) {Browser};
  \node[box,minimum width=3cm,right=4.4cm of b2] (a2) {API};
  \node[box,minimum width=2.4cm,right=3cm of a2] (r2) {Redis};
  \draw[arr] ([yshift=2mm]b2.east) -- node[lbl,above]{2. email + code} ([yshift=2mm]a2.west);
  \draw[arr] ([yshift=2mm]a2.east) -- node[lbl,above]{compare, delete} ([yshift=2mm]r2.west);
  \draw[arr] ([yshift=-2mm]a2.west) -- node[lbl,below]{access token + HttpOnly refresh cookie} ([yshift=-2mm]b2.east);
\end{tikzpicture}
\caption{Password, then one-time code}\label{fig:login}
\end{figure}

\subsection{Step 1: \code{POST \codesl{}api\codesl{}auth\codesl{}login\codesl{}request-otp}}

\begin{lstlisting}[language=JavaScript,caption={requestOtp (\file{src/modules/auth/auth.controller.js})}]
async function requestOtp(req, res) {
  const { email, password } = req.body;
  const user = await User.findOne({ email });

  if (!user) throw unauthorized(INVALID_CREDENTIALS, 'INVALID_CREDENTIALS');
  if (user.isAccountLocked()) {
    throw forbidden('Account is locked after too many failed attempts. Try again later.', 'ACCOUNT_LOCKED');
  }
  if (!user.isActive) throw forbidden('Account is deactivated. Contact the hospital administrator.', 'ACCOUNT_DEACTIVATED');

  if (!(await user.comparePassword(password))) {
    const failed = await recordFailedAttempt(user);
    return res.status(401).json({
      success: false,
      message: INVALID_CREDENTIALS,
      code: 'INVALID_CREDENTIALS',
      attemptsRemaining: Math.max(0, config.security.maxFailedLogins - failed),
    });
  }

  if (await otpService.getAttempts(email) >= config.otp.maxAttempts) {
    throw tooMany('Too many OTP requests. Try again in 15 minutes.');
  }

  const otp = otpService.generateOTP();
  await otpService.storeOTP(email, otp);
  await emailService.sendOTP(email, otp, user.fullName);
  await otpService.incrementAttempts(email);

  res.json({
    success: true,
    message: 'A login code was sent to your email',
    data: { email, expiresInMinutes: config.otp.expiryMinutes },
  });
}
\end{lstlisting}

\begin{howto}
\begin{enumerate}
  \item An unknown email and a wrong password give the same message
    (\code{INVALID\_CREDENTIALS}), so the form cannot be used to find out who
    has an account.
  \item A wrong password increments \code{failedLoginAttempts}; at
    \code{MAX\_FAILED\_LOGINS} (5) the account is locked for 30 minutes.
  \item A per-email counter in Redis limits how many codes can be requested.
  \item \code{sendOTP} is awaited and \emph{throws} on failure: if the email
    could not be sent the user is told (500), instead of waiting for a code that
    will never arrive. (A test checks exactly this.)
\end{enumerate}
\end{howto}

\begin{lstlisting}[language=JavaScript,caption={Lockout bookkeeping (\file{src/modules/auth/auth.controller.js})}]
async function recordFailedAttempt(user) {
  const updated = await User.findByIdAndUpdate(user._id, { $inc: { failedLoginAttempts: 1 } }, { new: true });
  if (updated.failedLoginAttempts >= config.security.maxFailedLogins) {
    updated.accountLockedUntil = new Date(Date.now() + config.security.lockMinutes * 60 * 1000);
    await updated.save({ validateBeforeSave: false });
    logger.warn(`Account locked after ${updated.failedLoginAttempts} failed attempts: ${user.email}`);
  }
  await invalidateSession(user._id); // trust score reflects the failure immediately
  return updated.failedLoginAttempts;
}
\end{lstlisting}

\code{\$inc} is atomic in MongoDB, so two simultaneous wrong guesses are both
counted. Recent failures also lower the trust score of any session still open
for that user (Section~\ref{sec:trust}).

\subsection{The OTP service: \file{src\codesl{}services\codesl{}otpService\codedot{}js}}

\begin{lstlisting}[language=JavaScript,caption={Generating, storing and checking codes (\file{src/services/otpService.js})}]
  /** Only the hash is stored: a Redis dump cannot be used to log in. */
  hashOTP(otp) {
    return crypto.createHash('sha256').update(String(otp)).digest('hex');
  }

  /** Cryptographically secure, uniformly random code with no leading zero. */
  generateOTP(length = config.otp.length) {
    return crypto.randomInt(10 ** (length - 1), 10 ** length).toString();
  }

  /** Stores the hashed code with a TTL; a new request overwrites the old code. */
  storeOTP(email, otp) {
    return guard('store', () => redisClient.setEx(otpKey(email), config.otp.expiryMinutes * 60, this.hashOTP(otp)));
  }

  /** Single use: a matching code is deleted immediately. */
  verifyOTP(email, otp) {
    return guard('verify', async () => {
      const stored = await redisClient.get(otpKey(email));
      if (!stored) return { success: false, message: 'OTP expired or not found' };

      const expected = Buffer.from(stored, 'hex');
      const actual = Buffer.from(this.hashOTP(otp), 'hex');
      // Constant-time comparison: response time does not leak how many characters matched.
      if (expected.length !== actual.length || !crypto.timingSafeEqual(expected, actual)) {
        return { success: false, message: 'Invalid OTP' };
      }
      await redisClient.del(otpKey(email));
      return { success: true, message: 'OTP verified successfully' };
    });
  }
\end{lstlisting}

\begin{howto}
\begin{itemize}
  \item \code{crypto.randomInt} is a cryptographically secure generator;
    \code{Math.random} is predictable and must never be used for codes.
  \item \code{setEx} stores the code with a time-to-live: Redis deletes it after
    10 minutes without any clean-up job.
  \item Only the SHA-256 of the code is stored, like a password.
  \item \code{timingSafeEqual} always takes the same time, so an attacker cannot
    learn from response timing how close a guess was.
  \item \code{guard} converts any Redis error into \code{REDIS\_UNAVAILABLE},
    which becomes a 503: if the OTP store is down, nobody logs in.
\end{itemize}
\end{howto}

\subsection{Step 2: \code{POST \codesl{}api\codesl{}auth\codesl{}login\codesl{}verify-otp}}

\begin{lstlisting}[language=JavaScript,caption={verifyOtp, success path (\file{src/modules/auth/auth.controller.js})}]
  const result = await otpService.verifyOTP(email, otp);
  if (!result.success) {
    const otpAttempts = await otpService.incrementAttempts(email);
    const failed = await recordFailedAttempt(user);
    return res.status(401).json({ /* ... INVALID_OTP, attemptsRemaining */ });
  }

  await otpService.clearAttempts(email);
  user.failedLoginAttempts = 0;
  user.accountLockedUntil = null;
  user.lastLogin = new Date();
  user.isVerified = true;
  await user.save({ validateBeforeSave: false });
  await invalidateSession(user._id);

  const access = await issueSession(res, user);
\end{lstlisting}

A correct code resets the failure counters and marks the account
\code{isVerified} (the user has proven they control the email address, which
raises their trust score). \code{issueSession} creates the token pair.

\section{Tokens: \file{src\codesl{}services\codesl{}tokenService\codedot{}js}}

MediChain uses two tokens with different jobs:

\begin{table}[H]
\centering\small
\caption{The two tokens}\label{tab:tokens}
\begin{tabularx}{\textwidth}{@{}P{2.6cm}YY@{}}
\toprule
 & \textbf{Access token} & \textbf{Refresh token} \\
\midrule
Purpose & Sent with every API call & Gets a new access token \\
Lifetime & 15 minutes & 7 days \\
Where the browser keeps it & A JavaScript variable (memory only) & HttpOnly cookie, path \code{\codesl{}api\codesl{}auth} \\
Signed with & \code{JWT\_SECRET} & \code{JWT\_REFRESH\_SECRET} \\
Revocation & Blacklist of its \code{jti} in Redis & Valid only while its Redis key exists \\
\bottomrule
\end{tabularx}
\end{table}

\begin{lstlisting}[language=JavaScript,caption={Issuing tokens (\file{src/services/tokenService.js})}]
  /** Short-lived access token. The random jti lets us revoke this exact token. */
  generateAccessToken(user) {
    const jti = crypto.randomUUID();
    const token = jwt.sign(
      { userId: user._id.toString(), email: user.email, role: user.role, jti },
      this.config.jwt.secret,
      { expiresIn: this.config.jwt.expiry }
    );
    return { token, jti, expiresIn: this.config.jwt.expiry };
  }

  /**
   * Long-lived refresh token, signed with a different secret. It is valid only
   * while its Redis key exists, so the server can revoke it at any time.
   */
  async generateRefreshToken(user) {
    const userId = user._id.toString();
    const jti = crypto.randomUUID();
    const token = jwt.sign(
      { userId, jti, type: 'refresh' },
      this.config.jwt.refreshSecret,
      { expiresIn: `${this.config.jwt.refreshExpiryDays}d` }
    );
    try {
      await this.redis.setEx(this.refreshKey(userId, jti), this.refreshTTLSeconds, '1');
    } catch (error) {
      logger.error(`Error storing refresh token: ${error.message}`);
      throw new Error(REDIS_UNAVAILABLE);
    }
    return { token, jti, jtiHash: this.hashJTI(jti) };
  }
\end{lstlisting}

\begin{lstlisting}[language=JavaScript,caption={The refresh cookie (\file{src/modules/auth/auth.controller.js})}]
const cookieOptions = () => ({
  httpOnly: true, // JavaScript in the page cannot read it (XSS)
  secure: config.isProduction, // HTTPS-only in production; plain http://localhost works in dev
  sameSite: 'strict', // never sent on cross-site requests (CSRF)
  path: '/api/auth', // only sent to the auth endpoints that need it
  maxAge: config.jwt.refreshExpiryDays * 24 * 60 * 60 * 1000,
});
\end{lstlisting}

\begin{good}[Why the token is not in localStorage]
Anything in \code{localStorage} can be read by any script on the page; one XSS
bug would leak every user's long-lived session. Here the long-lived token is in
an \code{HttpOnly} cookie that scripts cannot read, and the short-lived access
token exists only in a variable that disappears when the tab closes. A page
reload restores the session by calling \code{\codesl{}api\codesl{}auth\codesl{}refresh-token}, which the
browser authenticates with the cookie automatically.
\end{good}

\section{Refresh-token rotation}

\begin{lstlisting}[language=JavaScript,caption={refreshToken (\file{src/modules/auth/auth.controller.js})}]
async function refreshToken(req, res) {
  const presented = req.cookies?.[REFRESH_COOKIE];
  if (!presented) throw unauthorized('No session. Please log in.', 'NO_REFRESH_TOKEN');

  let decoded;
  try {
    decoded = await tokenService.verifyRefreshToken(presented);
  } catch (error) {
    if (error.message === 'REDIS_UNAVAILABLE') throw error;
    res.clearCookie(REFRESH_COOKIE, { ...cookieOptions(), maxAge: undefined });
    throw unauthorized('Session expired. Please log in again.', 'INVALID_REFRESH_TOKEN');
  }

  const user = await User.findById(decoded.userId);
  if (!user || !user.isActive) {
    await tokenService.revokeRefreshToken(decoded.userId, decoded.jti);
    throw unauthorized('Account deactivated', 'ACCOUNT_DEACTIVATED');
  }

  await tokenService.revokeRefreshToken(decoded.userId, decoded.jti);
  await User.updateOne({ _id: user._id }, { $pull: { refreshTokens: tokenService.hashJTI(decoded.jti) } });
  const access = await issueSession(res, user);
  // ... respond with the new access token and user
}
\end{lstlisting}

\begin{howto}
Each refresh token works \emph{once}. Using it deletes its Redis key and issues
a new one in a new cookie. If an attacker steals a refresh token and uses it,
the real user's next refresh fails (its token is gone), so the theft becomes
visible. If the real user refreshes first, the stolen copy is already dead.
\code{verifyRefreshToken} also checks \code{type === 'refresh'}, so an access
token cannot be replayed here.
\end{howto}

\section{Logout, password change and deactivation}

\begin{lstlisting}[language=JavaScript,caption={Revocation on logout (\file{src/modules/auth/auth.controller.js}) and blacklisting (\file{tokenService.js})}]
async function logout(req, res) {
  await tokenService.blacklistAccessToken(req.user);
  // ... verify and revoke the refresh cookie, clear it
}

  /** Blacklists an access token until it would have expired anyway. */
  async blacklistAccessToken(decodedPayload) {
    if (!decodedPayload?.jti || !decodedPayload?.exp) throw new Error('INVALID_TOKEN');
    const ttl = decodedPayload.exp - Math.floor(Date.now() / 1000);
    if (ttl <= 0) return;
    try {
      await this.redis.setEx(`jwt:blacklist:${decodedPayload.jti}`, ttl, '1');
    } catch (error) {
      // ...
      throw new Error(REDIS_UNAVAILABLE);
    }
  }
\end{lstlisting}

A JWT is valid until it expires; the server cannot ``un-sign'' it. So logout
puts the token's id on a blacklist that \code{verifyAccessToken} checks on every
request. The blacklist entry lives exactly as long as the token would have, so
the list never grows without bound.

Changing the password calls \code{revokeAllRefreshTokens}, which scans Redis for
every \code{refresh:<userId>:*} key and deletes them: every other device is signed
out. Deactivating a user in the admin dashboard does the same, plus
\code{invalidateSession} so the next request with a still-valid access token is
refused (Chapter~\ref{ch:zerotrust}).

\begin{warn}[A bug that was fixed]
The original controller was a class whose methods were passed to Express as
\code{router.post('\codesl{}logout', authController.logout)}. Express calls them without
an object, so \code{this} was \code{undefined} and logout always crashed. The
controller is now a set of plain functions with no \code{this}.
\end{warn}

\section{Redis key space}\label{sec:redis-keys}

\begin{table}[H]
\centering\small
\caption{What lives in Redis}\label{tab:redis}
\begin{tabularx}{\textwidth}{@{}P{5.2cm}P{2cm}Y@{}}
\toprule
\textbf{Key} & \textbf{TTL} & \textbf{Meaning} \\
\midrule
\code{otp:<email>} & 10 min & SHA-256 of the current login code \\
\code{otp:attempts:<email>} & 15 min & OTP requests + wrong codes \\
\code{refresh:<userId>:<sha256(jti)>} & 7 days & This refresh token is still valid \\
\code{jwt:blacklist:<jti>} & token's remaining life & This access token was revoked \\
\code{session:active:<userId>} & 5 min & Cached \code{isActive}, \code{isVerified}, \code{approvalStatus}, failures \\
\bottomrule
\end{tabularx}
\end{table}
